\subsection{随机逃逸与对噪声敏感的切换}

真实系统是带噪声的。
疲劳、中断、新奇性与情绪波动，都可以表现为状态扰动。
这一点很重要，因为跨越势垒既可能来自有意控制，也可能来自随机逃逸。

\begin{discussion}
随机视角有助于避免过度道德化。
一个人从良好的工作状态中漂移出来，并不一定是因为他积极选择了失败，
而可能只是因为这个状态过浅，且对噪声过于敏感。
反过来，一个好的干预之所以有效，也未必是因为它提升了“德性”，
而可能只是因为它把目标盆地加深到了足以让普通噪声不再轻易把轨迹击出该区域。
\end{discussion}